\subsection{First Theorems on Peano Systems}

\begin{topicbox}{Induction as a Working Principle}
\begin{exposition}
The induction axiom for a Peano system is formulated in terms of inductive
subsets of the underlying set. Arithmetic arguments are usually written in
terms of predicates rather than subsets. This theorem identifies the
predicate-form induction principle derived from the Peano axioms. It is the
standard form of induction used in the rest of the chapter.
\end{exposition}

\begin{tcolorbox}[colback=thmbox, colframe=thmborder, arc=2pt,
  left=6pt, right=6pt, top=4pt, bottom=4pt,
  title={\small\textbf{Theorem (Induction Principle for a Peano System)}},
  fonttitle=\small\bfseries]
\begin{theorem}[Induction Principle for a Peano System]
\label{thm:induction-principle-for-peano-system}
Let $(P,S,1)$ be a Peano system, and let $\varphi$ be a predicate on $P$.
If
\[
\varphi(1)
\qquad\text{and}\qquad
\forall n \in P\;(\varphi(n) \Longrightarrow \varphi(S(n))),
\]
then
\[
\forall n \in P\;\varphi(n).
\]
\hyperref[prf:induction-principle-for-peano-system]{Go to proof.}
\end{theorem}
\end{tcolorbox}

\begin{remark*}[Standard quantified statement]
\[
\begin{aligned}
&\text{Let } (P,S,1) \text{ be a Peano system and let } \varphi
\text{ be a predicate on } P.\\
&\Bigl(\varphi(1) \land
\forall n \in P\;(\varphi(n) \Longrightarrow \varphi(S(n)))\Bigr)
\Longrightarrow
\forall n \in P\;\varphi(n).
\end{aligned}
\]
\end{remark*}

\begin{remark*}[Contrapositive quantified statement]
\[
\begin{aligned}
&\text{Let } (P,S,1) \text{ be a Peano system and let } \varphi
\text{ be a predicate on } P.\\
&\neg\forall n \in P\;\varphi(n)
\Longrightarrow
\neg\Bigl(\varphi(1) \land
\forall n \in P\;(\varphi(n) \Longrightarrow \varphi(S(n)))\Bigr).
\end{aligned}
\]
\end{remark*}

\begin{remark*}[Interpretation]
This theorem converts the subset form of induction into the predicate form
used in ordinary arithmetic arguments. A property propagates through the
entire Peano system once it holds at the distinguished first element and is
preserved by successor. The theorem therefore turns the structural induction
axiom into a proof principle for arbitrary predicates on $P$.
\end{remark*}

\begin{remark*}[Comparison with Feferman]
Feferman derives the working form of mathematical induction from the subset
formulation of the Peano axioms immediately after Theorem 3.3 in Chapter 3 of
\citet{FefermanNumberSystems1964}. The present notes formalize that induction
principle as an explicit reusable theorem.
\end{remark*}

\begin{remark*}[Dependencies]
\begin{itemize}
  \item \hyperref[def:peano-system]{Peano System}
  \item \hyperref[ax:peano-induction]{Induction Axiom}
\end{itemize}
\end{remark*}
\end{topicbox}

\begin{topicbox}{Inductive Subsets and Predecessors}
\begin{exposition}
Induction is expressed in terms of subsets of a Peano system. Two structural
notions appear repeatedly in that setting: closure under successor and
membership generated from the distinguished first element. Predecessor language
is also convenient once it is known that every non-initial element arises as a
successor.
\end{exposition}

\begin{tcolorbox}[colback=defbox, colframe=defborder, arc=2pt,
  left=6pt, right=6pt, top=4pt, bottom=4pt,
  title={\small\textbf{Definition (Successor-Closed Subset)}} ,
  fonttitle=\small\bfseries]
\begin{definition}[Successor-Closed Subset]
\label{def:successor-closed-subset}
Let $(P,S,1)$ be a Peano system, and let $A \subseteq P$. The subset $A$ is
called \emph{successor-closed} exactly when
\[
\forall n \in P\;(n \in A \Longrightarrow S(n) \in A).
\]
\end{definition}
\end{tcolorbox}

\begin{remark*}[Standard quantified statement]
\[
\begin{aligned}
&\text{Let } (P,S,1) \text{ be a Peano system and let } A \subseteq P.\\
&A \text{ is successor-closed}
\Longleftrightarrow
\forall n \in P\;(n \in A \Longrightarrow S(n) \in A).
\end{aligned}
\]
\end{remark*}

\begin{remark*}[Negated quantified statement]
\[
\begin{aligned}
&\text{Let } (P,S,1) \text{ be a Peano system and let } A \subseteq P.\\
&A \text{ is not successor-closed}
\Longleftrightarrow
\exists n \in P\;(n \in A \land S(n) \notin A).
\end{aligned}
\]
\end{remark*}

\begin{remark*}[Failure modes]
A subset fails to be successor-closed exactly when it contains some element of
$P$ but fails to contain that element's successor.
\end{remark*}

\begin{remark*}[Failure mode decomposition]
\[
\begin{aligned}
A \text{ is not successor-closed}
\Longleftrightarrow
\exists n \in P\;\underbrace{(n \in A \land S(n) \notin A)}_{\text{closure fails at } n}.
\end{aligned}
\]
\end{remark*}

\begin{remark*}[Interpretation]
A successor-closed subset is stable under one step of counting. Once an
element of $A$ is present, its successor remains in $A$ as well.
\end{remark*}

\begin{remark*}[Comparison with Feferman]
The notion of a successor-closed subset is implicit in condition (iii) of
Feferman's Definition 3.1 \citep{FefermanNumberSystems1964}. The present notes
isolate it as a separate definition so that induction arguments can refer to
the closure condition directly.
\end{remark*}

\begin{remark*}[Dependencies]
\begin{itemize}
  \item \hyperref[def:peano-system]{Peano System}
\end{itemize}
\end{remark*}

\begin{tcolorbox}[colback=defbox, colframe=defborder, arc=2pt,
  left=6pt, right=6pt, top=4pt, bottom=4pt,
  title={\small\textbf{Definition (Inductive Subset of a Peano System)}} ,
  fonttitle=\small\bfseries]
\begin{definition}[Inductive Subset of a Peano System]
\label{def:inductive-subset-of-peano-system}
Let $(P,S,1)$ be a Peano system, and let $A \subseteq P$. The subset $A$ is
called \emph{inductive} exactly when $1 \in A$ and $A$ is successor-closed.
\end{definition}
\end{tcolorbox}

\begin{remark*}[Standard quantified statement]
\[
\begin{aligned}
&\text{Let } (P,S,1) \text{ be a Peano system and let } A \subseteq P.\\
&A \text{ is inductive}
\Longleftrightarrow
\bigl(1 \in A \land
\forall n \in P\;(n \in A \Longrightarrow S(n) \in A)\bigr).
\end{aligned}
\]
\end{remark*}

\begin{remark*}[Negated quantified statement]
\[
\begin{aligned}
&\text{Let } (P,S,1) \text{ be a Peano system and let } A \subseteq P.\\
&A \text{ is not inductive}
\Longleftrightarrow
\bigl(1 \notin A \lor \exists n \in P\;(n \in A \land S(n) \notin A)\bigr).
\end{aligned}
\]
\end{remark*}

\begin{remark*}[Failure modes]
A subset fails to be inductive in exactly two ways: it may fail to contain the
distinguished first element, or it may contain some element while failing to
contain that element's successor.
\end{remark*}

\begin{remark*}[Failure mode decomposition]
\[
\begin{aligned}
A \text{ is not inductive}
\Longleftrightarrow\;&
\underbrace{1 \notin A}_{\text{base point missing}}\\
&\lor
\underbrace{\exists n \in P\;(n \in A \land S(n) \notin A)}_{\text{successor closure fails}}.
\end{aligned}
\]
\end{remark*}

\begin{remark*}[Interpretation]
An inductive subset contains the distinguished first element and is closed
under successor. Such a subset contains every stage obtained by starting at
$1$ and applying successor repeatedly.
\end{remark*}

\begin{remark*}[Comparison with Feferman]
The notion of an inductive subset packages the two hypotheses appearing in the
induction clause of Feferman's Definition 3.1
\citep{FefermanNumberSystems1964}. The present notes name this package so
that the internal structure of induction proofs remains explicit.
\end{remark*}

\begin{remark*}[Dependencies]
\begin{itemize}
  \item \hyperref[def:peano-system]{Peano System}
  \item \hyperref[def:successor-closed-subset]{Successor-Closed Subset}
\end{itemize}
\end{remark*}

\begin{tcolorbox}[colback=defbox, colframe=defborder, arc=2pt,
  left=6pt, right=6pt, top=4pt, bottom=4pt,
  title={\small\textbf{Definition (Predecessor in a Peano System)}} ,
  fonttitle=\small\bfseries]
\begin{definition}[Predecessor in a Peano System]
\label{def:predecessor-in-peano-system}
Let $(P,S,1)$ be a Peano system, and let $x,u \in P$. The element $u$ is
called a \emph{predecessor} of $x$ exactly when
\[
S(u)=x.
\]
\end{definition}
\end{tcolorbox}

\begin{remark*}[Standard quantified statement]
\[
\begin{aligned}
&\text{Let } (P,S,1) \text{ be a Peano system and let } x,u \in P.\\
&u \text{ is a predecessor of } x
\Longleftrightarrow
S(u)=x.
\end{aligned}
\]
\end{remark*}

\begin{remark*}[Negated quantified statement]
\[
\begin{aligned}
&\text{Let } (P,S,1) \text{ be a Peano system and let } x,u \in P.\\
&u \text{ is not a predecessor of } x
\Longleftrightarrow
S(u) \neq x.
\end{aligned}
\]
\end{remark*}

\begin{remark*}[Interpretation]
A predecessor of $x$ is an element whose successor is $x$. This language
isolates the inverse-looking question attached to the successor map without
assuming that successor is globally invertible.
\end{remark*}

\begin{remark*}[Dependencies]
\begin{itemize}
  \item \hyperref[def:peano-system]{Peano System}
\end{itemize}
\end{remark*}
\end{topicbox}

\begin{topicbox}{Generation by Successor}
\begin{exposition}
The Peano axioms describe a counting system generated from a distinguished
first element by the successor operation. The next theorems make that
generation principle explicit: every element is either the first element or a
successor, every non-initial element has a unique predecessor, and no element
can be its own successor.
\end{exposition}

\begin{tcolorbox}[colback=thmbox, colframe=thmborder, arc=2pt,
  left=6pt, right=6pt, top=4pt, bottom=4pt,
  title={\small\textbf{Theorem (Every Element Is Either \texorpdfstring{$1$}{1} or a Successor)}} ,
  fonttitle=\small\bfseries]
\begin{theorem}[Every Element Is Either $1$ or a Successor]
\label{thm:every-element-is-one-or-a-successor}
Let $(P,S,1)$ be a Peano system. For every $x \in P$, either $x=1$ or there
exists $u \in P$ such that $S(u)=x$.
\hyperref[prf:every-element-is-one-or-a-successor]{Go to proof.}
\end{theorem}
\end{tcolorbox}

\begin{remark*}[Standard quantified statement]
\[
\begin{aligned}
&\text{Let } (P,S,1) \text{ be a Peano system.}\\
&\forall x \in P\;\bigl(x=1 \lor \exists u \in P\;(S(u)=x)\bigr).
\end{aligned}
\]
\end{remark*}

\begin{remark*}[Interpretation]
Every element of a Peano system is reached by starting at the distinguished
first element and moving forward by successor. The only possible exception to
having a predecessor is the element $1$ itself.
\end{remark*}

\begin{remark*}[Historical note]
This theorem corresponds directly to Theorem 3.3 in
\citet{FefermanNumberSystems1964}. Feferman uses this result to show that the
distinguished initial element is the only possible non-successor in a Peano
system.
\end{remark*}

\begin{remark*}[Dependencies]
\begin{itemize}
  \item \hyperref[def:peano-system]{Peano System}
  \item \hyperref[ax:peano-one-in-n]{Distinguished Element Lies in the Underlying Set}
  \item \hyperref[ax:peano-successor-closure]{Closure Under Successor}
  \item \hyperref[thm:induction-principle-for-peano-system]{Induction Principle for a Peano System}
\end{itemize}
\end{remark*}

\begin{tcolorbox}[colback=corbox, colframe=corborder, arc=2pt,
  left=6pt, right=6pt, top=4pt, bottom=4pt,
  title={\small\textbf{Corollary (Predecessor Existence and Uniqueness Away from \texorpdfstring{$1$}{1})}} ,
  fonttitle=\small\bfseries]
\begin{corollary}[Predecessor Existence and Uniqueness Away from $1$]
\label{cor:predecessor-exists-unique-away-from-one}
Let $(P,S,1)$ be a Peano system, and let $x \in P$ with $x \neq 1$. Then
there exists a unique $u \in P$ such that $S(u)=x$.
\hyperref[prf:predecessor-exists-unique-away-from-one]{Go to proof.}
\end{corollary}
\end{tcolorbox}

\begin{remark*}[Standard quantified statement]
\[
\begin{aligned}
&\text{Let } (P,S,1) \text{ be a Peano system.}\\
&\forall x \in P\;\bigl(x \neq 1 \Longrightarrow \exists! u \in P\;(S(u)=x)\bigr).
\end{aligned}
\]
\end{remark*}

\begin{remark*}[Interpretation]
Every element distinct from $1$ has exactly one predecessor. The successor map
therefore organizes the non-initial part of a Peano system into one-step
extensions of earlier elements.
\end{remark*}

\begin{remark*}[Comparison with Feferman]
Feferman proves predecessor existence and uniqueness together inside Theorem
3.3 of \citet{FefermanNumberSystems1964}. The present notes separate the
uniqueness statement into an independent corollary so that it can be cited
directly in later proofs.
\end{remark*}

\begin{remark*}[Dependencies]
\begin{itemize}
  \item \hyperref[def:peano-system]{Peano System}
  \item \hyperref[thm:every-element-is-one-or-a-successor]{Every Element Is Either $1$ or a Successor}
  \item \hyperref[ax:peano-successor-injective]{Injectivity of Successor}
\end{itemize}
\end{remark*}

\begin{tcolorbox}[colback=corbox, colframe=corborder, arc=2pt,
  left=6pt, right=6pt, top=4pt, bottom=4pt,
  title={\small\textbf{Corollary (The Distinguished Element Is the Unique Non-Successor)}} ,
  fonttitle=\small\bfseries]
\begin{corollary}[The Distinguished Element Is the Unique Non-Successor]
\label{cor:one-is-unique-non-successor}
Let $(P,S,1)$ be a Peano system. An element $x \in P$ is not a successor of
any element of $P$ if and only if $x=1$.
\hyperref[prf:one-is-unique-non-successor]{Go to proof.}
\end{corollary}
\end{tcolorbox}

\begin{remark*}[Standard quantified statement]
\[
\begin{aligned}
&\text{Let } (P,S,1) \text{ be a Peano system.}\\
&\forall x \in P\;\Bigl(\neg\exists u \in P\;(S(u)=x) \Longleftrightarrow x=1\Bigr).
\end{aligned}
\]
\end{remark*}

\begin{remark*}[Interpretation]
The distinguished first element is characterized internally as the unique
element with no predecessor. This gives a structural description of $1$ that
does not depend on external notation.
\end{remark*}

\begin{remark*}[Comparison with Feferman]
This corollary extracts a structural characterization of the distinguished
initial element that is implicit in the discussion surrounding Theorem 3.3 of
\citet{FefermanNumberSystems1964}.
\end{remark*}

\begin{remark*}[Dependencies]
\begin{itemize}
  \item \hyperref[ax:peano-one-not-successor]{Distinguished Element Is Not a Successor}
  \item \hyperref[cor:predecessor-exists-unique-away-from-one]{Predecessor Existence and Uniqueness Away from $1$}
\end{itemize}
\end{remark*}

\begin{tcolorbox}[colback=thmbox, colframe=thmborder, arc=2pt,
  left=6pt, right=6pt, top=4pt, bottom=4pt,
  title={\small\textbf{Theorem (No Object Is Its Own Successor)}} ,
  fonttitle=\small\bfseries]
\begin{theorem}[No Object Is Its Own Successor]
\label{thm:no-object-is-its-own-successor}
Let $(P,S,1)$ be a Peano system. For every $n \in P$,
\[
S(n) \neq n.
\]
\hyperref[prf:no-object-is-its-own-successor]{Go to proof.}
\end{theorem}
\end{tcolorbox}

\begin{remark*}[Standard quantified statement]
\[
\begin{aligned}
&\text{Let } (P,S,1) \text{ be a Peano system.}\\
&\forall n \in P\;(S(n) \neq n).
\end{aligned}
\]
\end{remark*}

\begin{remark*}[Interpretation]
Successor always moves forward in the counting structure. No element can be
recovered by applying successor to itself.
\end{remark*}

\begin{remark*}[Historical note]
This theorem is a standard consequence of the Peano axioms in the classical
development of the natural numbers. Feferman includes the corresponding result
as an exercise in Chapter 3 of \citet{FefermanNumberSystems1964}.
\end{remark*}

\begin{remark*}[Dependencies]
\begin{itemize}
  \item \hyperref[ax:peano-one-not-successor]{Distinguished Element Is Not a Successor}
  \item \hyperref[ax:peano-successor-injective]{Injectivity of Successor}
  \item \hyperref[thm:induction-principle-for-peano-system]{Induction Principle for a Peano System}
\end{itemize}
\end{remark*}
\end{topicbox}

